% ============================== Part G ==============================
\problem{G1}{Bernstein \& Armstrong for a constant shear}

\begin{given}
BA14 write the posterior of one shear shared by many galaxies as $P(\vg\mid D) \propto P(\vg)\prod_iP(D_i\mid\vg)$. They
expand each galaxy's likelihood to second order about $\vg = 0$,
$P(D_i\mid\vg) \approx P_i + \vg\cdot\mathbf Q_i + \tfrac12\vg\cdot\mathbf R_i\cdot\vg$ (here $\mathbf R_i$ is BA14's
second-derivative matrix, not our rotation), and find a Gaussian posterior with
$C_g^{-1} = \sum_i(\mathbf Q_i\mathbf Q_i^\T/P_i^2 - \mathbf R_i/P_i)$ and $\bar\vg = C_g\sum_i\mathbf Q_i/P_i$.
\source{\cBAconst.}
\end{given}

\need{A3, A4, B3, B4.}

\begin{enumerate}
  \item Show that $\mathbf Q_i/P_i = \nabla\ln P(D_i\mid\vg)|_0$ and
    $\mathbf Q_i\mathbf Q_i^\T/P_i^2 - \mathbf R_i/P_i = -\nabla\nabla\ln P(D_i\mid\vg)|_0 \equiv \mathbf A_i$.
  \item For a Gaussian likelihood $\ln L_i = -\tfrac12(\vg - \hat\vg_i)^\T F_i(\vg - \hat\vg_i)$, show that $\mathbf A_i = F_i$
    and $\nabla\ln L_i|_0 = F_i\hat\vg_i$, and hence that BA14's $\bar\vg$ is our stack $\Ftot^{-1}\sum_iF_i\hat\vg_i$ (A3).
  \item Our likelihood is $L_i = q_i/\pitrain$ (A4) on a $401^2$ grid, multiplied without Taylor expansion. What do we gain
    over the quadratic form, and what does the Gaussian form still do for us?
\end{enumerate}

\solution{
(a) $\nabla\ln P = \nabla P/P$, which at 0 is $\mathbf Q_i/P_i$. $\nabla\nabla\ln P = \nabla\nabla P/P - \nabla P\nabla P^\T/P^2$,
which at 0 is $\mathbf R_i/P_i - \mathbf Q_i\mathbf Q_i^\T/P_i^2$.
(b) $\nabla\ln L_i = -F_i(\vg - \hat\vg_i)$, which at 0 is $F_i\hat\vg_i$, and $-\nabla\nabla\ln L_i = F_i$. So
$C_g^{-1} = \sum_iF_i$ and $\bar\vg = (\sum_iF_i)^{-1}\sum_iF_i\hat\vg_i$.
(c) The full product is exact at any $|g| < 0.1$ and includes the higher-order terms BA14 add as perturbations
(\cBAthird). The Gaussian form still defines the reference shear: BA14's $\mathbf A_i$ is our $F_i$, so $\gref$ uses
the same weights the stack does (B4).
}

% --------------------------------------------------------------------
\problem{G2}{Bernstein \& Armstrong for a varying shear}

\begin{given}
For any model with parameters $\mathbf a$ that predicts a shear $\vg_i(\mathbf a)$ at each galaxy, BA14's quadratic
expansion gives (written per galaxy)
\[
-\ln P(\mathbf a\mid D) \approx \text{const} - \ln P(\mathbf a) - \sum_i\vg_i(\mathbf a)\cdot\frac{\mathbf Q_i}{P_i}
+ \frac12\sum_i\vg_i(\mathbf a)\cdot\mathbf A_i\cdot\vg_i(\mathbf a).
\]
If $\vg_i$ is linear in $\mathbf a$ and the prior is Gaussian, the maximum has a closed form.
\source{\cBAvarying.}
\end{given}

\need{G1, B4, E2, E4.}

\begin{enumerate}
  \item \emph{Cells.} Let $\mathbf a$ be one shear per cell, with $\vg_i(\mathbf a) = R_i^\T\vG_{c(i)}$. Show that the
    posterior separates into independent cell stacks. Using $\E[\mathbf Q_i/P_i] \approx \mathbf A_i\vt_i$ (the linear
    response of B4), show that the stack converges to $\gref$.
  \item \emph{$\kappa$ map.} Let $\mathbf a$ be the $\kappa$ pixels, with $\vg_i(\kappa)$ from Part D. Neglecting the
    $1/(1-\kappa)$ factor, $\vg_i$ is linear in $\kappa$: write $\vg = M\kappa$. With the prior of E4, which closed form
    does BA14's statement give, and what is it called?
  \item \emph{Halo fit.} What is $\mathbf a$, and is the model linear?
  \item In one sentence: what does a galaxy contribute to any shear-field model?
\end{enumerate}

\solution{
(a) Each galaxy's term depends only on its own cell's $\vG_c$, so the sum splits by cell, giving one A3/G1 stack per
cell. Setting the gradient with respect to $\vG_c$ to zero, $\sum_{i\in c}\big(R_i\mathbf A_iR_i^\T\vG_c - R_i\mathbf Q_i/P_i\big) = 0$,
so $\hat\vG_c = (\sum_iF_i')^{-1}\sum_iR_i\mathbf Q_i/P_i$ with $F_i' = R_i\mathbf A_iR_i^\T$. With
$\E[\mathbf Q_i/P_i] = \mathbf A_i\vg_i = \mathbf A_iR_i^\T\vt_i$, each term becomes $F_i'\vt_i$, and
$\E[\hat\vG_c] = (\sum_iF_i')^{-1}\sum_iF_i'\vt_i = \gref$.
(b) Write the linear forward model as $\vg = M\kappa$ ($M$ stacks the rows $\partial\vg_i/\partial\kappa$), let
$\mathcal A = \mathrm{blockdiag}(\mathbf A_i)$ and $\mathbf q$ stack the $\mathbf Q_i/P_i$, and take the Gaussian prior
$S^{-1}$ of E4. The objective is quadratic, so $\hat\kappa = (S^{-1} + M^\T\mathcal AM)^{-1}M^\T\mathbf q$. This is E2's
estimator with $\mathcal A$ in the role of $N^{-1}$ (and $\mathbf q = \mathcal A\hat\vg$ for Gaussian likelihoods): the
Wiener filter. Our map is its non-linear, non-Gaussian generalisation, solved numerically.
(c) $\mathbf a = (M_{\rm main}, M_{\rm sub})$; the NFW shear is non-linear in the masses, so it needs sampling or
optimisation.
(d) Its log-likelihood evaluated at the shear the model predicts at that galaxy's position: no per-galaxy ``shear
estimate'' is ever needed.
}

% --------------------------------------------------------------------
\problem{G3}{Why the difference map is clean}

\begin{given}
In each realization, the main-only and main + sub datasets have the same galaxies and the same noise draws; only the
true shears differ. The map estimator is (nearly) the linear MAP of E2, $\hat\kappa = \mathsf H^{-1}A^\T N^{-1}\vd$. The
difference map is $\hat\kappa_{\rm main+sub} - \hat\kappa_{\rm main}$. \emph{Derived here}
(\texttt{KAPPA\_FITTING\_EXPLAINED.md}, Sec.~9.4).
\end{given}

\need{E2, F3, F4.}

\begin{enumerate}
  \item In the linear model, show that the difference map equals $\mathsf R\,\kappa_{\rm subhalo}$ \emph{exactly}, with
    no noise at all.
  \item Our real difference maps have a background rms of $\approx0.002$ against a peak of $\approx0.014$. Where can
    noise come from, if (a) says there is none?
  \item Why can't a real observation make a difference map, and what does it show instead?
  \item The subhalo's true peak pixel is 0.030 and the difference map peaks at about 0.012--0.014. Explain with
    $\mathsf R$.
\end{enumerate}

\solution{
(a) $\vd_{\rm sub} - \vd_{\rm main} = A(\kappa_{\rm sub} - \kappa_{\rm main}) = A\kappa_{\rm subhalo}$, because the noise is
identical and cancels. Then the difference map is $\mathsf H^{-1}A^\T N^{-1}A\kappa_{\rm subhalo} = \mathsf R\,\kappa_{\rm subhalo}$.
(b) From the non-linearity: $g = \gamma/(1-\kappa)$ couples the subhalo to the main halo, and the NPE likelihoods are not
exactly Gaussian, so the noise does not cancel perfectly.
(c) A real cluster gives one map; there is no noise-matched twin without the subhalo. The difference map shows how
visible the subhalo is \emph{through the method}: its position, size and amplitude after the prior's smoothing. The
significance a real map could claim comes from G4.
(d) The rows of $\mathsf R$ near the subhalo are broad and low (F4(d)): the map shrinks and spreads it to about 40\% of
its peak.
}

% --------------------------------------------------------------------
\problem{G4}{\py\ The significance one map can claim}

\begin{given}
For each of $K$ independent realizations of the cluster (new galaxies, positions and noise), we average the
reconstructed $\kappa$ over a 0.5\,Mpc disk at the subhalo position: $a_{{\rm sub},r}$ for the main + sub map,
$a_{{\rm main},r}$ for the main-only map. The single-map significance is
$z_r = (a_{{\rm sub},r} - \overline{a_{\rm main}})/{\rm sd}(a_{\rm main})$, with the mean and sd over realizations. An
aperture average of a mass map is a standard detection statistic (\cSchneiderMap); the empirical-null $z$ is
\emph{derived here}.
\end{given}

\need{B1, G3.}

\begin{enumerate}
  \item Show that subtracting the same fixed map (for example the true main-halo $\kappa$) from every map leaves every
    $z_r$ unchanged.
  \item Why is $z_r$ computed against the other realizations' main-only maps, and not against realization $r$'s own
    main-only map?
  \item Our null is built from maps of the \emph{true} main halo. For a real cluster it would come from simulations of
    a \emph{fitted} main halo. Why does that make our $z$ an upper bound?
  \item With our numbers, $\overline{a_{\rm main}} = 0.0072$, ${\rm sd} = 0.0038$ and $\overline{a_{\rm sub}} = 0.0156$,
    compute the $z$ of the average realization.
  \item \py\ In the 1D toy of E5, generate 50 realizations with and without a narrow bump, compute $z_r$, and estimate
    the fraction with $z > 2$.
\end{enumerate}

\solution{
(a) Subtracting a fixed map subtracts the same constant $c$ (its aperture average) from every $a_{{\rm sub},r}$ and every
$a_{{\rm main},r}$. The difference in the numerator and the sd are unchanged.
(b) Realization $r$'s own main-only map shares its noise, which would cancel (G3), and a real observation has no such
twin. The other realizations represent ``what that spot looks like in a cluster without a subhalo'', across noise and
galaxy layouts.
(c) Uncertainty in the fitted mass, concentration and centre adds scatter to the null, widening its sd and lowering $z$.
(d) $(0.0156 - 0.0072)/0.0038 = 2.2$ (the median over realizations is 2.4).
(e) The notebook (one seed) gives median $z \approx 1.9$, with 48\% of realizations above $z = 2$: a toy version of our 2.4 and 60\%. It also confirms (a): subtracting the truth changes no $z$.
}

% --------------------------------------------------------------------
\problem{G5}{Reading the ensemble numbers}

\begin{given}
Our results over 50 independent NPE realizations (\texttt{repo\_reports/PAPER\_DRAFT\_HANDOFF.md}): one realization
constrains $m$ to $\approx\pm0.3$; the pooled ensemble gives $m = +0.04 \pm 0.04$ (polar $g_t$); the pull spread is 0.84;
the likelihood map's median $z$ is 2.4 (KS: 1.1); and the best possible matched-filter S/N (C5(d), with full precision
matrices) has median 3.7 over realizations. Statistical errors average down as $1/\sqrt K$ (A2).
\end{given}

\need{A2, B1, C5.}

\begin{enumerate}
  \item Check $0.3/\sqrt{50}$ against the ensemble error. How many realizations would $\pm0.01$ need? Why doesn't the
    cluster test aim for that?
  \item What fraction of the matched-filter S/N does the likelihood map achieve, and what fraction does KS achieve?
    Name two reasons the map falls short of the matched filter.
  \item With 4 times the source density (Xu sample 1), roughly what median $z$ would you expect, if $z \propto \sqrt{n}$?
\end{enumerate}

\solution{
(a) $0.3/\sqrt{50} = 0.042$, consistent. $\pm0.01$ needs $(0.3/0.01)^2 = 900$ realizations, because one cluster's cells
span only $\approx0.04$ in shear and the slope is poorly pinned. The test asks a different question: no extra bias
from shear gradients at the few-percent level. Percent-level $m$ comes from the constant-shear stacks.
(b) $2.4/3.7 = 65\%$; KS $1.1/3.7 = 30\%$. The matched filter knows the subhalo's exact profile and position, while the
map assumes neither. The map's prior also shrinks the subhalo and the null includes galaxy-layout variance.
(c) $\approx2.4\times2 = 4.8$.
}

% --------------------------------------------------------------------
\problem{G6}{Why the map settings were chosen on synthetic data}

\begin{given}
We tried 27 likelihood-map settings and 4 KS settings. The rule ``highest contrast with a false-peak rate $\le 0.10$''
was fixed before the NPE maps were looked at, and applied to synthetic data only. For a standard normal, $P(z > 2) =
0.0228$. \emph{Derived here}.
\end{given}

\need{B1, G4.}

\begin{enumerate}
  \item If instead we had tried all 31 settings on the real map and reported the best $z$, and the settings were
    independent, what is the chance that at least one gives $z > 2$ when there is no subhalo?
  \item Real settings are correlated. Does that make the problem better or worse than (a)?
  \item Why does choosing on synthetic data, with the rule fixed in advance, avoid the problem?
\end{enumerate}

\solution{
(a) $1 - (1 - 0.0228)^{31} = 0.51$: a coin flip.
(b) Better, since correlated settings behave like fewer independent tries, but the bias stays positive.
(c) The choice never sees the data it is judged on, so the reported $z$ is a single, pre-registered test.
}
